\begin{lemma}
\label{lemma-constructible-stable-specialization-closed}
Let $R$ be a ring. Let $E \subset \Spec(R)$ be a constructible subset.
\begin{enumerate}
\item If $E$ is stable under specialization, then $E$ is closed.
\item If $E$ is stable under generalization, then $E$ is open.
\end{enumerate}
\end{lemma}

\begin{proof}
First proof. By Lemma \ref{lemma-constructible}, retain
a finite expression
$$
E=\bigcup_{i=1}^a\left(D(f_i)\cap
V(g_{i1},\ldots,g_{ib_i})\right).
$$
For $J_i=(g_{i1},\ldots,g_{ib_i})$, the original
quotient and localization maps give
$\Spec((R/J_i)_{f_i})$ the image $D(f_i)\cap V(J_i)$.
Contraction is the map; the inverse on this subset
sends $\mathfrak p$ to $(\mathfrak p/J_i)_{f_i}$.
The finite product $A=\prod_i(R/J_i)_{f_i}$ therefore
has spectrum image $E$ under the diagonal original
map $R\to A$, as in Lemma \ref{lemma-constructible-is-image}.
Indeed the coordinate idempotents sum to $1$ and
have pairwise zero products, so every prime selects
exactly one factor; this identifies its contraction
with that factor's contraction. The empty product
is the zero ring and gives the empty set.
If $E$ is stable under specialization, the preceding
image lemma proves it closed.

The complement $E^c$ is constructible by
Topology, Lemma \ref{topology-lemma-constructible}.
If $E$ is stable under generalization, $E^c$ is
stable under specialization: a specialization from
a point outside $E$ cannot enter $E$, since the
reverse generalization would then enter $E$ too.
Applying the first assertion to $E^c$ proves it
closed, hence $E$ open.
Each original basic term $D(f_i)\cap V(J_i)$ is
quasi-compact, being the image of the spectrum of
$(R/J_i)_{f_i}$. Therefore every constructible subset
of $\Spec(R)$ is quasi-compact by this finite union.
The complement is constructible as well. Thus in (1)
the closed set has quasi-compact complement, and in
(2) the open set is quasi-compact.

\medskip\noindent
Second proof, including its stronger scope. By
Lemma \ref{lemma-spec-spectral}, $X=\Spec(R)$ is
spectral. Its constructible topology has the
quasi-compact opens and their complements as a
subbasis; these sets are therefore both open and
closed in that topology. The topology is
quasi-compact Hausdorff by Topology, Lemma
\ref{topology-lemma-constructible-hausdorff-quasi-compact}.
Let $C\subseteq X$ be closed in the constructible
topology, which is weaker as a hypothesis than
being constructible, and let $x\in\overline C$
for the original topology. For every quasi-compact
open neighbourhood $U$ of $x$, the set $U\cap C$
is nonempty and is closed in the constructible
topology. Finite intersections of these $U$ are
again quasi-compact open neighbourhoods, so this
family has the finite intersection property.
Compactness gives $y\in C$ lying in every such $U$.
Since these sets form a neighbourhood basis of $x$,
$x$ is a specialization of $y$. Thus
$$
\overline C=\{x:\text{some }y\in C\text{ specializes to }x\}.
$$
If $C$ is stable under specialization it is closed
in the original topology. By taking complements,
any subset open in the constructible topology
and stable under generalization is open in the
original topology. This gives the full stronger
form of Topology, Lemma
\ref{topology-lemma-constructible-stable-specialization-closed},
as well as both original assertions. The stronger
scope does not assert that these sets or their
complements are constructible or quasi-compact in
all the cases just stated.
\end{proof}